\section*{Kurzfassung}
Die Betrag/Phase-Zerlegung des FFGFT-Massenoperators (Dok.~288/289) ordnet die Informationsfrage neu.
Information hat eine \emph{Betragsseite} (Masse/Energie -- Landauer und Vopson, von FFGFT aus der
$T^4$-Geometrie hergeleitet) und eine unabhängige \emph{Phasenseite} (Holonomie, Mischung -- das eine
offene Datum $\theta$). Der reinste Beleg ist das \textbf{Photon}: exakt masselos -- null Ruhemasse-Betrag, im Massenkanal also reine Phase -- und das
Quantum des $U(1)$-Eichpotentials $A_\mu$, das die Verbindung ist; FFGFTs Informationseinheit $E_\text{bit}=\hbar c/L$ \emph{ist}
eine Photonenenergie. Die \emph{Menge} an Information ist dabei nicht darstellungs-relativ, sondern
absolut und dimensionslos in $\xi$ verankert: $I_1=\log_2(1/\xi)\approx12{,}87$ Bit pro Freiheitsgrad;
Energie und Kapazität sind ihre dimensionalen, frequenz-skalierenden Beiwerke. Die Frage ist nicht
\enquote{ist Information fundamental?}, sondern \textbf{welcher Kanal} -- und nur die Geometrie hält
beide. Strukturelle Diagnose, keine neue Herleitung.

\section{Ausgangslage: Geometrie zuerst, und die Betrag/Phase-Brille}
FFGFTs Position zur Informationsfrage ist dokumentiert und geometrie-zuerst. Dok.~243 (Diskussion mit
Onur Teker, UIFT): die Abbildung von FFGFT in einen Informationsformalismus ist eine \emph{Projektion}
$\pi$, kein Isomorphismus; die Bit-Zählung $I=\log_2|\mathcal{W}|$ \emph{setzt} den geometrischen
Zustandsraum \emph{voraus}. Dok.~255/257 leiten Landauers Prinzip und Vopsons Masse-Information-Äquivalenz
\emph{aus} der $T^4$-Geometrie ab: ein Bit ist eine minimale topologische Windung $\Delta w=1$,
\enquote{Information ist topologisch eingefrorene kinetische Energie}, $E_\text{bit}(L)=\hbar c/L=pc$.
Dok.~274 stellt FFGFT \emph{geometrie-primitiv} neben informations-primitiv (Vopson). FFGFT lehnt Vopson
also nicht ab -- es begründet ihn.

Die neue Brille ist die Betrag/Phase-Zerlegung des Massenoperators (Dok.~288/289): der Operator ist ein
Z$_3$-Zirkulant, DFT-diagonalisiert, und sein Inhalt zerfällt in ein \emph{Betragsspektrum} (Massen,
Energie) und ein \emph{Phasenspektrum} (Mischung, Holonomie). Der $\theta$-Befund (Dok.~282) sagt: das
Betragsspektrum legt das Phasenspektrum nur bis auf einen Allpass-Faktor fest (präzisiert in §2). Das ist der Schlüssel zu allem Folgenden.

\section{Zwei Kanäle: Betrag und Phase}
$\Delta w=1$ und $E_\text{bit}=\hbar c/L$ sind Betrags-/Energie-Aussagen. Da das Betragsspektrum das
Phasenspektrum nur bis auf einen Allpass-Faktor festlegt (siehe unten), folgt:

Vopsons Masse-Information-Äquivalenz -- die FFGFT herleitet -- erfasst strukturell \emph{nur den
Betragskanal}. Der Phasenkanal (Holonomie, Mischung, $\theta$) ist unabhängige Information, nicht auf
Masse/Energie reduzierbar.

Das zeigt sich an zwei Stellen konkret. \emph{Erstens, Mischungsblindheit:} ein Massenoperator
$M=U\,\mathrm{diag}(m)\,U^{\dagger}$ trägt die Eigenwerte $m$ (Betragskanal) und die Eigenvektoren $U$
(Mischung, Phasenkanal); jede Funktion der Eigenwerte allein ($\sum m$, Koide $Q$) ist gegenüber $U$
invariant. Numerisch (\texttt{masse\_information\_phasenblind.py}): bei festen Massen $(0{,}5,1{,}5,3{,}0)$
bleiben $\sum m=5$ und $Q=0{,}372$ über drei Mischungen gleich, während $\lVert M_\text{offdiag}\rVert$
($1{,}42$ bis $1{,}78$) variiert -- zwei physikalisch verschiedene Operatoren mit derselben
Masse-Information.

\emph{Zweitens, die Phase trägt die Struktur:} nachrichtentechnisch seit Oppenheim/Lim (1981) bekannt --
tauscht man Betrags- und Phasenspektrum zweier Muster über Kreuz, folgt die Rekonstruktion dem
\emph{Phasen}-Spender. Numerisch (\texttt{phase\_traegt\_struktur.py}): Phase-only korreliert mit dem
Original ($+0{,}53$), Magnitude-only nicht ($-0{,}24$). Eine Masse-Energie-Buchhaltung erfasst also den
\emph{strukturärmeren} Kanal.

\emph{Drittens, präzise:} \enquote{der Betrag legt die Phase nicht fest} gilt nur \emph{bis auf einen
Allpass-Faktor}. Jede stabile Transferfunktion zerfällt in $H=H_\text{min}\cdot H_\text{ap}$: für den
\emph{minimalphasigen} Anteil legt der Betrag die Phase sehr wohl fest (Bode-/Hilbert-Relation,
$\arg=\mathcal{H}\{\log|H|\}$), während der \emph{Allpass} ($|H_\text{ap}|=1$) eine Phase trägt, die der
Betrag \emph{nicht} festlegt. Numerisch (\texttt{minimalphase\_allpass.py}): die allein aus dem Betrag
rekonstruierte Phase trifft $\arg H_\text{min}$ (Fehler $\sim10^{-15}$), die Restphase ist exakt der
Allpass. Frei -- und damit das Offene -- bleibt der \emph{Allpass-Anteil}, und genau dort sitzt die
$\theta$-Holonomie (Lesart Minimalphase\,(FFGFT)\,$\times$\,Allpass\,(HLV-$R$), Dok.~288).

\section{Das Photon: der exakte Anker}
Der reinste Beleg ist das Photon -- und der Korpus enthält die Teile bereits. Dok.~007 behandelt
Neutrinos über die Photonen-Analogie: Photon $E^2=(pc)^2+0$ (exakt masselos), Neutrino
$E^2=(pc)^2+(\sqrt{\xi^2/2}\,mc^2)^2$ (quasi-masselos). Das Photon ist der $m=0$-Grenzfall, das Neutrino
\enquote{Photon plus erste $\xi^2/2$-Korrektur}.

Das Photon hat Ruhemasse-Betrag exakt null -- der exakte Allpass \emph{im Massenkanal}: der gesamte Ruhemasse-Inhalt ist Phase. (Die Energie $\hbar\omega$ bleibt eine Betragsgröße; \enquote{reine Phase} meint den Ruhemasse-, nicht den Energiekanal.) Drei Dinge
fallen zusammen:
\begin{enumerate}
  \item \textbf{Ungefrorene Information.} Das Photon ist der ungefrorene Fall von
        \enquote{Information $=$ eingefrorene kinetische Energie} ($E=pc$); FFGFTs Einheit
        $E_\text{bit}=\hbar c/L$ \emph{ist} eine Photonenenergie, Landauers thermischer Bit ihre
        Einfrierung.
  \item \textbf{Photon und Verbindung.} Das Eichpotential $A_\mu$ \emph{ist} die $U(1)$-Verbindung; das
        Photon ist deren Quantum. \emph{Ein} eichinvarianter Observable ist die Holonomie
        ($\oint A\,\mathrm{d}l$, Aharonov-Bohm, sogar wo $F=0$) -- neben Energie $\hbar\omega$, Impuls und
        Polarisation, die keine Holonomien sind. Der Bezug zu \enquote{it from connection} ist
        \emph{begrifflich} (relationaler Charakter), keine formale Identität mit $\theta$.
  \item \textbf{Gegenbeispiel zu \enquote{Masse $=$ Information}.} Null Ruhemasse, doch Träger praktisch
        unserer gesamten Informationsübertragung.
\end{enumerate}

Das Neutrino ist nur die erste Korrektur daran: ein \enquote{gedämpftes Photon} mit winzigem Betrag statt
exakt null. Im entarteten Grenzfall wäre $M=m\cdot\mathbb{1}$ für jedes $U$
($\lVert M-m\mathbb{1}\rVert\sim10^{-16}$), die Mischung verließe $M$, der Oszillationsinhalt säße in der
Propagations-/Holonomie-Phase -- dieselbe Lage, die das Photon \emph{exakt} und ohne
Entartungs-Hypothese realisiert.

\section{Was \enquote{Informationsmenge} heißt: Paket, Zahl, Energie, Kapazität}
Ist ein Photon das kleinste Paket? Nur \emph{pro Mode}. Ein Einzelfrequenz-Photon ist eine ebene Welle;
ein lokalisiertes Photon ist ein Wellenpaket. Wichtig: ein zeitlich begrenztes Paket
behält den \emph{definiten} Träger $f_0$ (die Laserfarbe, der Energiesprung); die Begrenzung erzeugt nur
ein Band $\Delta f\sim1/\Delta t$ \emph{um} $f_0$ (Fensterwirkung), kein Spektrum. Die Zyklenzahl ist
$N=f_0\,\Delta t$: bei festem Zeitfenster $\propto f_0$, aber bei fester Zyklenzahl ist die fraktionale
Bandbreite $\Delta f/f_0=1/N$ \emph{frequenz-unabhängig}. Beim Photon aus einem Energiesprung ist die
Lebensdauer $\tau$ das Fenster: die natürliche Linienbreite $\Delta f=f_0/Q$ mit $Q=\omega_0\tau=2\pi N$
($\tau=10$\,ns, $f_0=5\cdot10^{14}$\,Hz $\Rightarrow N=5\cdot10^6$, $\Delta f\approx16$\,MHz; numerisch
\texttt{fenster\_zyklen\_linienbreite.py}). Definiter Träger \emph{plus} Fenster ergibt das Band.

\paragraph{Weder Periodenzahl noch Fensterlänge ist intrinsisch.} Beide setzt der Emissionsprozess, nicht
das \enquote{Photon-Sein}: eine monochromatische Mode hat unendlich viele Perioden und unendliches Fenster
(ebene Welle); ein spontan emittiertes Photon erbt $f_0$ \emph{und} $\tau$ vom Übergang
($N=f_0\tau=Q/2\pi$, beides gekoppelt); ein technisches Paket fixiert eine Größe, die andere folgt aus
$N=f_0\Delta t$. Intrinsisch -- skalen-invariant -- ist die dimensionslose Zahl $N$ ($\Delta f/f=1/N$),
$\Delta t=N/f$ ist ihre dimensionale Einkleidung. Die Spannweite ist groß: Few-Cycle-Puls $N\approx2$,
schmale Atomlinie $N\approx5\cdot10^6$ (numerisch \texttt{quant\_bit\_modeninfo.py}) -- es gibt keine
universelle Periodenzahl und keine universelle Fensterlänge.

Damit sind drei Größen sauber zu trennen (numerisch \texttt{photon\_paket\_count\_energie.py}):

\begin{itemize}
  \item \textbf{Zahl}: die Quantenzahl (Besetzung) ist ein \emph{Zählwert} -- eine Anregung,
        frequenz-unabhängig, relational --, \emph{kein} Bitmaß. Was ein Quant an \emph{Information} trägt,
        steht unten.
  \item \textbf{Energie} pro Bit: $E=\hbar\omega=\hbar c/L$ -- frequenz-abhängig (Radio
        $\sim2\cdot10^{-10}$\,eV bis Gamma $\sim2\cdot10^{5}$\,eV).
  \item \textbf{Kapazität}: Shannon $C=B\log_2(1+\mathrm{SNR})$, $B\propto f$ -- frequenz-abhängig
        ($\sim2\cdot10^{5}$ bis $\sim2\cdot10^{20}$ Bit/s).
\end{itemize}

\enquote{Höhere Frequenz $\to$ mehr Information} ist genau die Kapazität: mehr Bandbreite, mehr Moden pro
Sekunde. Das ist eine \emph{andere} Achse als Betrag/Phase: \emph{innerhalb} einer Mode trägt jedes
Symbol Amplitude \emph{und} Phase (I/Q bzw.\ QAM), \emph{über} die Moden zählt die Bandbreite. FFGFT
spiegelt beides: innerhalb der Mode der Z$_3$-Zirkulant ($r,\theta$), über die Moden $|\mathcal{W}|$.

\paragraph{Ein Quant ist ein Zählwert, kein Bit.} \enquote{Ein Quant $=$ ein Bit} ist keine allgemeine
Gleichung -- sie verwechselt Zählen mit Messen. Ein Quant ist eine \emph{Anregung} (Zählwert~$1$); sein
\emph{Informationsgehalt} ist $\log_2 M$ Bit, wenn es aus $M$ unterscheidbaren Moden gewählt ist --
\emph{variabel}, pro Freiheitsgrad gedeckelt bei $I_1=\log_2(1/\xi)\approx12{,}9$ Bit (§5). \enquote{Ein
Quant $=$ ein Bit} gilt nur für $M=2$ (Mode auf ein Binär-Alphabet fixiert) -- eine Kodierkonvention, kein
Naturgesetz. Zählwert und Bitgehalt sind verschiedene Achsen: ein Quant in $2^{13}$ Moden trägt $13$~Bit,
$13$ binäre Quanten ebenso -- gleiche Bitzahl, verschiedene Quantenzahl (numerisch
\texttt{quant\_bit\_modeninfo.py}). Das Quant \emph{zählt}; $\log_2 M\le I_1$ \emph{misst}. Damit ist auch
§5 sauber eingeordnet: $I_1$ ist nicht der Gehalt \emph{eines} Quants, sondern die absolute \emph{Decke}
pro Freiheitsgrad, die der variable Modengehalt $\log_2 M$ nie überschreitet.

\section{Die absolute Menge: in $\xi$ verankert}
Eine frühere Lesart -- der Informationsgehalt sei darstellungs-relativ, es gebe keine absolute Bit-Zahl
-- war überzogen und FFGFT-untreu. Der absolute, frequenz-unabhängige Gehalt ist die \emph{dimensionslose}
Größe, und $\xi$ verankert sie. $\xi=4/30000$ ist dimensionslos (der absolute Parameter); $\hbar,c$ sind
SI-Buchhaltung.

$\xi$ deckelt die Windungszahl bei $n_\text{max}\sim1/\xi$ (Dok.~243/009) -- ohne diese Schranke
divergierte $I=\log_2|\mathcal{W}|$. Eine Windungskoordinate kann also $\sim1/\xi$ unterscheidbare Werte
annehmen: $I_1=\log_2(1/\xi)=\log_2(7500)\approx12{,}87$~Bit. Gesamt
$I_\text{gesamt}\sim N_\text{max}\cdot\log_2(1/\xi)$ (Zellen $N_\text{max}=V/L_0^{d}$, Dok.~251, mal
Tiefe). Dimensionslos, von $\xi$ gesetzt -- die Energie $\hbar c/L$ ist nicht die Menge, sondern ihr
dimensionaler Preis, und die Kapazität sättigt an der $\xi$-Skala.

Nur \emph{schwach} relativ bleibt die Bit-Zahl einer \emph{Beschreibung} bei gewählter Auflösung und das,
was eine verlustbehaftete Projektion verwirft (Dok.~243) -- eine Aussage über Näherungen, nicht über die
Information selbst. Die volle Information (Geometrie bzw.\ ihr verlustfreies Hilbert-Bild, Dok.~230) ist
absolut.

\paragraph{Absolut und relational sind kein Widerspruch.} Die beiden Aussagen -- die Menge ist
\emph{absolut} (§5), das offene Datum $\theta$ ist \emph{relational} (§8) -- liegen auf verschiedenen
Achsen: \enquote{absolut} misst die \emph{Menge}/das Maß, \enquote{relational} beschreibt die
\emph{Natur} eines Datums. Eine Holonomie (Windung) ist das Musterbeispiel für \emph{beides}: relational
definiert (Schleifenintegral, kein lokaler Punktwert) \emph{und} absolut bewertet (eichinvariant,
ganzzahlig). Numerisch (\texttt{absolut\_relational.py}): eine Windung $w=3$ ist nur als Schleifenintegral
ablesbar (die halbe Schleife liefert $1$) und bleibt unter einwertiger Umphasung unverändert
(eichinvariant); und die absolute Zahl $\log_2(1/\xi)$ \emph{zählt} genau solche relationalen Objekte
($\sim15000$ Windungssektoren bis $n_\text{max}\sim1/\xi$, $\log_2\approx13{,}9$ Bit). Die Absolutheit
sitzt im Boden (der Geometrie, $\xi$), die Relationalität in den Freiheitsgraden darauf -- die absolute
Bit-Zahl zählt relationale Windungen.

\section{Statische und dynamische Definition des Bits}
Die bisherige Definition ist durchweg \emph{statisch}: das Bit als Struktur -- Windung $\Delta w=1$,
Zähldecke $\log_2 M\le I_1$, Energie $E_0=\hbar c/L$ (Betragsseite). Das sagt, \emph{was} ein Bit ist,
nicht was es \emph{tut}. Die dynamische Definition ist kein Zusatz, sondern die konjugierte Seite -- und
sie zerfällt selbst in zwei, in denen die Energie \emph{verschiedene} Rollen spielt.

\emph{Kinematisch (reversibel), die $\tilde T$-Seite:} ein Flip dauert $\tilde T=\hbar/E_0=1/m_0$, die Rate
ist $1/\tilde T=E_0/\hbar$. Hier ist $E_0\cdot\tilde T=\hbar$ nichts anderes als $\tilde T\cdot m=1$ -- die
Zeit-Masse-Dualität, das FFGFT-Fundament. Das Quanten-Speed-Limit (Margolus-Levitin) $\tau\ge\pi\hbar/2E$
macht es scharf: die Energie \emph{setzt die Uhr}, ohne dissipiert zu werden (unitär). Die Kosten sind
\emph{Zeit}. Die Kapazitätsachse (Bit/s, §4) ist genau diese Rate.

\emph{Thermodynamisch (irreversibel), die Landauer-Seite:} ein Bit trägt Entropie $k_B\ln 2$; das
\emph{Löschen} (Vergessen) in ein Bad bei Temperatur $T$ dissipiert $k_B T\ln 2$ Wärme. Das hängt an $T$,
\emph{nicht} an $E_0$, und die Wärme wird \emph{dissipiert} -- die Kosten sind \emph{Wärme}, nicht Zeit.
In FFGFT ist diese Seite \emph{emergent}: das fundamentale $T^4$-Geschehen ist unitär-reversibel, sodass
kinematisch ($\tilde T$) primär ist und der Landauer-Preis erst durch Grobkörnung, Bad und logische
Irreversibilität entsteht.

Der dimensionslose Zählwert $\log_2 M$ ist primär; die Konstanten kleiden ihn ein (numerisch
\texttt{bit\_statisch\_dynamisch.py}): $\hbar,c$ zu Energie $E_0$ und Zeit $\tilde T$ (kinematisch,
reversibel: \emph{Zeit}-Kosten, $\tilde T\cdot m=1$), $k_B$ zu Entropie $k_B\ln 2$ und Löschwärme
$k_B T\ln 2$ (thermodynamisch, irreversibel: \emph{Wärme}-Kosten). $\hbar,c,k_B$ sind SI-Buchhaltung.
Die Energie spielt zwei Rollen: \emph{Uhr} (kinematisch, kein Verlust) versus \emph{dissipierte Wärme}
(thermodynamisch).

Wichtig: \enquote{Buchhaltung} heißt nicht \emph{entbehrlich}. $k_B,\hbar,c$ sind die \emph{notwendigen}
Einheiten-Brücken. SI trennt Temperatur (K), Energie (J), Zeit (s) in eigene Einheiten, und sobald
thermodynamisch \emph{gerechnet} wird -- und das geschieht in SI --, ist $k_B$ die unverzichtbare
Umrechnung K$\to$J (sonst stünde die Löschwärme in Kelvin statt Joule), wie $\hbar,c$ die Umrechnung zu
$E_0=\hbar c/L$ tragen. In natürlichen Einheiten ($k_B=\hbar=c=1$) sind die Kosten reine Zahlen ($\ln 2$);
SI führt die Einheitenwände ein, und die Buchhaltung ist deren Mörtel -- ontologisch sekundär, operativ
unverzichtbar (numerisch \texttt{buchhaltung\_si\_natuerlich.py}).

\paragraph{Das Bit in SI-Zahlen.} Verankert an der charakteristischen Energie des Rahmens,
$E_0=\sqrt{m_e m_\mu}=7{,}348$\,MeV (über $\alpha=\xi E_0^2$ an die Feinstrukturkonstante gekoppelt; die
Variante $\sqrt{\alpha/\xi}=7{,}398$\,MeV weicht um $+0{,}68\%$ ab -- keine exakte Gleichheit). Das
generische $E_0=\hbar c/L$ fällt mit diesem $E_0$ genau bei $L=L_0=\hbar c/E_0$ zusammen. Dann (numerisch
\texttt{bit\_in\_si.py}): $E_0=1{,}177\cdot10^{-12}$\,J, $L_0=2{,}686\cdot10^{-14}$\,m $=26{,}85$\,fm,
$m_0=1{,}310\cdot10^{-29}$\,kg, $\tilde T=\hbar/E_0=8{,}958\cdot10^{-23}$\,s $=L_0/c$ und
$1/\tilde T=1{,}116\cdot10^{22}$\,s$^{-1}$.

\section{Die Temperatur des Bits: Skala, Zustand, Bad}
Die Energie $E_0$ lässt sich direkt in eine \emph{Temperatur} umrechnen. Über $E_0=k_B T_0$ ist
\[
T_0=\frac{E_0}{k_B}=\frac{1{,}177\cdot10^{-12}\,\text{J}}{1{,}381\cdot10^{-23}\,\text{J/K}}
   =8{,}53\cdot10^{10}\,\text{K}
\]
(Probe über $1$\,eV$=11\,604{,}5$\,K: $7{,}348\cdot10^6\,\text{eV}\cdot11\,604{,}5=8{,}53\cdot10^{10}$\,K),
also rund $85$ Milliarden Kelvin. $T_0$ ist die Temperatur, bei der die thermische Energie $k_B T_0$ die
\emph{Struktur-Energie} $E_0$ des Bits erreicht.

\paragraph{Hat ein Bit überhaupt eine Temperatur?} Streng genommen: nein, nicht \emph{eine}. Ein einzelnes
Bit in definitem Zustand ($0$ oder $1$) hat Nullentropie und keine Verteilung -- Temperatur ist eine
Ensemble-/Reservoir-Größe und einem einzelnen, festgelegten Freiheitsgrad gar nicht zugeordnet. Was man
einem Bit sinnvoll zuordnen kann, sind \emph{drei verschiedene} temperaturartige Größen, die man nicht
verwechseln darf. Erstens die \emph{Zustands-} oder Besetzungstemperatur: hat das Bit zwei um $E_0$
getrennte Niveaus, so definiert das Besetzungsverhältnis $p_1/p_0=e^{-E_0/k_B T}$ eine effektive
Temperatur. Sie läuft \emph{gegen} die Anschauung: maximal gemischt ($p_0=p_1=\tfrac12$, volle ein Bit
Entropie) heißt $T_\text{eff}\to\infty$, definit ($p\to0$ oder $1$, null Entropie) heißt
$T_\text{eff}\to0$, Besetzungsumkehr sogar $T_\text{eff}<0$. Das \emph{informationsreichste} Bit ist damit
der Unendlich-Temperatur-Zustand; Kühlen treibt es in einen definiten Zustand und \emph{zerstört} die
Information. (Ein rein logisches, entartetes Bit -- beide Werte gleiche Energie -- hat gar keine
Besetzungstemperatur.) Zweitens die \emph{Skala} $T_0=E_0/k_B$: nicht \enquote{wo das Bit sitzt}, sondern
der charakteristische Schmelz-/Aktivierungspunkt aus dem Energie-Gap, die natürliche Temperatur-\emph{Einheit}
des Bits. Drittens die \emph{Bad}-Temperatur, die in die Löschkosten eingeht -- und die gehört der
\emph{Umgebung}, nicht dem Bit.

Die Skala $T_0$ markiert die thermische Belastbarkeitsgrenze: darüber ($k_B T>E_0$) überwiegt die
thermische Energie die Struktur, und die Wärme wäscht die Distinktion -- die Windung bzw.\ Mode, die das
Bit \emph{ist} -- aus: das Bit \enquote{schmilzt}. Darunter ($k_B T<E_0$) ist es robust. Das ist eine
Aussage über die \emph{Skala}, nicht über einen Zustand.

\paragraph{Welche Temperatur beim Löschen?} Die Landauer-Kosten $Q=k_B T\ln2$ hängen an der Temperatur
des \emph{Bades}, in das die gelöschte Entropie als Wärme abfließt -- also an der Wahl der Umgebung, nicht
am Bit. Es gibt keinen universellen \enquote{richtigen} Wert. Für einen realen Rechner ist es die
Betriebstemperatur, praktisch \emph{Raumtemperatur} ($\sim300$\,K): $Q=k_B\cdot300\cdot\ln2=18$\,meV. Das
ist die ehrliche Antwort für reales Rechnen. UIFT setzt statt dessen $T_\text{CMB}=2{,}7$\,K ein -- nicht
den absoluten Nullpunkt, aber die \emph{kälteste natürliche Wärmesenke}; das gibt die minimal mögliche
Löschwärme im heutigen Universum ($Q=0{,}16$\,meV, kälter geht ohne aktives Kühlen nicht). Legitim als
kosmologischer Boden, aber weiterhin nur \emph{eine} Bad-Wahl, keine Bit-Eigenschaft. $T_0$ schließlich
ist \emph{kein} Bad: setzt man es ein, bekommt man den Grenzfall \enquote{Bad so heiß wie das Bit selbst},
$Q=E_0\ln2=0{,}69\,E_0$ -- eine Referenz, kein Löschvorgang.

\enquote{Die} Temperatur eines Bits hängt damit an seiner \emph{Skala}: ein generisches Bit der Länge $L$
trägt die Energie $E=\hbar c/L$ und somit die Skala $T=\hbar c/(k_B L)$. Jede Bit-Skala hat ihre eigene;
$8{,}5\cdot10^{10}$\,K ist die des charakteristischen $E_0$. Ein größeres $L$ -- ein \enquote{weicheres}
Bit -- hat eine niedrigere Schwelle.

Damit steht $T_0$ neben Masse und Zeit als dritte Einheiten-Kleidung \emph{derselben} Energie. Mit je
einer Buchhaltungskonstante:
\[
m_0 c^2 \;=\; E_0 \;=\; k_B T_0 \;=\; \hbar/\tilde T,
\]
also $/c^2\to$ Masse $m_0=1{,}310\cdot10^{-29}$\,kg, $/\hbar\to$ Zeit $\tilde T=8{,}958\cdot10^{-23}$\,s
(die Dual-Zeit aus $\tilde T\cdot m=1$), $/k_B\to$ Temperatur $T_0=8{,}53\cdot10^{10}$\,K. Unter allen vier
sitzt der dimensionslose Zählwert $\log_2 M$. Eine Energie, vier Kleidungen.

Die Temperatur ist dabei nicht bloß \enquote{auch noch da}, sondern direkt mit der dynamischen Dual-Zeit
verknüpft. Die \emph{thermische Zeit} bei $T_0$ ist $\hbar/(k_B T_0)$; Einsetzen von $T_0=E_0/k_B$ gibt
$\hbar/E_0=\tilde T$ -- die thermische Zeit bei $T_0$ \emph{ist} die Dual-Zeit aus $\tilde T\cdot m=1$
(numerisch identisch bis $10^{-38}$\,s, \texttt{bit\_temperatur.py}). Das ist die KMS- bzw.\
\enquote{thermische Zeit}-Korrespondenz (Connes--Rovelli): Temperatur ist eine inverse (imaginäre) Zeit,
und sie landet ohne Zusatzannahme genau auf der FFGFT-Dual-Zeit. Auf der Buchhaltungsachse ist $k_B$ für
die Temperatur damit exakt das, was $\hbar$ für die Zeit ist.

Schließlich verbindet $T_0$ die beiden \enquote{Energien} eines Bits -- was es \emph{ist} ($E_0$) und was
sein \emph{Löschen} kostet ($k_B T\ln2$) -- zu \emph{einer} Skala. Die Landauer-Löschwärme relativ zur
Eigenenergie ist schlicht das Temperaturverhältnis,
\[
\frac{Q}{E_0}=\frac{k_B T\ln2}{k_B T_0}=\frac{T}{T_0}\,\ln 2.
\]
Bei Raumtemperatur: $\tfrac{300}{8{,}53\cdot10^{10}}\ln2=2{,}4\cdot10^{-9}$, also
$E_0/Q=(T_0/T)/\ln2\approx4\cdot10^{8}$ -- die Raumtemperatur liegt rund $8{,}5$ Dekaden unter $T_0$, und
\emph{daher} liegen die statische Eigenenergie ($\sim7{,}3$\,MeV) und die Löschwärme ($\sim18$\,meV) gut
acht Größenordnungen auseinander. Bei $T=T_0$ dagegen kostet das Löschen $E_0\ln2=0{,}69\,E_0$: am eigenen
Temperatur-Punkt sind \emph{Sein} und \emph{Löschen} dieselbe Größenordnung. Genau dieses Verhältnis
$T/T_0$ ist die \emph{einzige} aufbau-unabhängige Aussage: ein Bit hat keine absolute Temperatur, nur
seine Bad-Temperatur relativ zur Eigenskala. Statische ($E_0,T_0$) und thermodynamische ($k_B T\ln2$)
Seite sind so über das eine Verhältnis verbunden -- kein Sprung, eine Skala.

Einordnung: $T_0\approx8{,}5\cdot10^{10}$\,K ist die MeV-/Kernskala -- die Temperatur-Größenordnung der
Kern- und Hadronenphysik: Sternkerne (Sonnenkern $\sim1{,}5\cdot10^7$\,K) und das Quark-Gluon-Plasma in
Schwerionen-Stößen (QCD-Übergang $\sim1{,}7\cdot10^{12}$\,K bei $150$\,MeV). Die \enquote{heißes
Frühuniversum}-Lesart bleibt bewusst außen vor -- sie ist eine $\Lambda$CDM-Annahme, die FFGFT nicht
teilt (FFGFTs CMB-Struktur folgt aus der T$^4$-Geometrie, nicht aus einer heißen Frühphase). Ehrlich:
$T_0=E_0/k_B$ ist die \emph{Standard}-Energie-Temperatur-Zuordnung; FFGFT-spezifisch ist allein, dass sie
über $\tilde T\cdot m=1$ mit der Dual-Zeit zusammenfällt -- eine saubere Konsistenz, keine neue Herleitung.
Und weil die absolute Löschtemperatur die des Bades ist, wäre ein an ein bestimmtes Bad ($T_\text{CMB}$
oder Raumtemperatur) gekoppeltes $\xi$ von dieser Wahl abhängig -- genau der Grund, warum das
CMB-gebundene $\xi_\text{UIFT}$ eine kontingente, andere Größe ist und FFGFTs geometrisches $\xi$
bad-\emph{un}abhängig bleibt (§9).

\paragraph{$E_0$ ist eine Sprosse, nicht der Boden.} Die Skala $E_0$ ist nicht das unterste Ende des
Rahmens, sondern eine \emph{Sprosse} einer fraktalen Leiter. Der geometrische Boden liegt bei der
Sub-Planck-Länge $\xi\,\ell_P\approx2{,}2\cdot10^{-39}$\,m (Dok.~163) -- der kleinsten Länge, unter der die
Raumzeit keine wohldefinierten Freiheitsgrade mehr trägt. Weil $\xi\ell_P$ um den Faktor $\xi$
\emph{unter} der Planck-Länge liegt, steht dieser Boden in den \enquote{harten} Kleidungen um $1/\xi=7500$
\emph{über} Planck:
\begin{center}
\begin{tabular}{@{}lll@{}}
Länge & $\xi\ell_P$ & $2{,}16\cdot10^{-39}$\,m\\
Zeit/Dauer & $\xi t_P$ & $7{,}19\cdot10^{-48}$\,s\\
Frequenz & $1/(\xi t_P)$ & $1{,}39\cdot10^{47}$\,Hz\\
Energie & $E_P/\xi$ & $9{,}16\cdot10^{22}$\,GeV\\
Temperatur & $T_P/\xi$ & $1{,}06\cdot10^{36}$\,K\\
Masse & $m_P/\xi$ & $1{,}63\cdot10^{-4}$\,kg\\
\end{tabular}
\end{center}
Dieser Boden ist der \emph{kleinste-Länge}-Rand und damit in Energie/Temperatur/Frequenz das
\emph{Maximum} (UV): \enquote{Boden} meint feinste Auflösung, nicht kleinste Energie. Die Bit-Sprosse
$E_0$ ($7{,}348$\,MeV, $26{,}85$\,fm, $T_0=8{,}5\cdot10^{10}$\,K) liegt rund $25$ Dekaden \emph{gröber}.
Dazwischen läuft das Fraktal nach innen, Faktor $1/\xi=7500$ pro Hauptstufe (Dok.~146), mit je
$\approx13$ $\sqrt2$-Zwischenschritten -- und $\log_2 7500\approx12{,}87$ ist \emph{genau} die
Informationsdecke $I_1$ pro Freiheitsgrad aus §5. Die Zahl der Fraktal-Zwischenschritte pro Stufe
\emph{ist} die maximale Bitzahl; Kapazität des Bits und Feinheit des Fraktals sind derselbe Zählwert.

\paragraph{Landauer.} Das Löschen eines Bits ist logisch irreversibel: zwei Zustände werden auf einen
abgebildet, die Entropie sinkt um $k_B\ln2$. Nach dem zweiten Hauptsatz muss diese Entropie als Wärme an
die Umgebung abfließen, $Q\ge k_B T\ln2$ (Landauer 1961). Dabei ist $T$ die Temperatur des \emph{Bades},
in das gelöscht wird -- nicht die des Bits.

\paragraph{Die Problematik.} Und hier klemmt es. Der Rahmen liefert jetzt eine ganze Leiter
\emph{definiter} Temperaturen -- die Sprosse $T_0=8{,}5\cdot10^{10}$\,K, den Boden
$T_{\max}=1{,}1\cdot10^{36}$\,K, dazwischen jede Fraktalstufe ihre eigene, über $25$ Dekaden. \emph{Keine}
davon ist \enquote{die} Landauer-Temperatur. Setzt man eine intrinsische Skala in $Q=k_B T\ln2$ ein,
bekommt man keinen realen Löschvorgang, sondern einen Grenzfall oder Unsinn: mit $T_0$ die Referenz
$Q=0{,}69\,E_0$, mit $T_{\max}$ eine absurd große Wärme, mit $T_\text{CMB}$ die UIFT-Zahl (§9). Der reale
Löschvorgang nimmt die Temperatur der \emph{tatsächlichen Umgebung} (Computer: $\sim300$\,K). Das ist die
Problematik: trotz einer Leiter definiter Eigen-Temperaturen bleibt die Löschtemperatur die des Bades, und
invariant ist allein das \emph{Verhältnis} $Q/E=(T/T_\text{Sprosse})\ln2$ -- Bad relativ zur jeweiligen
Sprosse. Ein Bit hat eine definite Energie und, je Sprosse, eine definite Skalen-Temperatur; einen
definiten \emph{Löschpreis} bekommt es erst mit einer Umgebung.

\section{Folgerungen}
\emph{Die treue Beschreibung ist quanten-, nicht bit-basiert.} Weil der Hilbertraum eine Bijektion ist
(Dok.~230), die klassische Bit-Zählung aber eine Projektion: die verlustfreie Beschreibung von FFGFT ist
die Amplitude-plus-Phase-Darstellung (Quanteninformation); der Shannon-Bit-Count ist ihr verlustbehafteter
Schatten. Ein Informations-Primitiv braucht Betrag \emph{und} Phase.

\emph{Das Offene ist relationale Information.} Die eine irreduzible offene Größe (Dok.~282) ist die
Holonomie-Phase $\theta$ -- wegabhängig, relational, kein lokaler Bit-Gehalt. Das Residuum, das die
Bit-Projektion verwirft, \emph{ist} die Phase; und die Phase \emph{ist}, wo die offene Physik liegt. Die
fundamentale Information sitzt in der Verbindung: das elektromagnetische Feld $A_\mu$ ist diese Verbindung,
das Photon ihr Quantum -- ein begriffliches, kein numerisch mit $\theta$ identisches Beispiel.
\enquote{It from bit} wird \enquote{it from connection}.

\section{Ehrliche Grenzen}
\textbf{Strukturelle Diagnose, keine neue Herleitung} (P35). Dok.~255/257 hatten Landauer/Vopson schon;
neu ist die Betrag/Phase-Zerlegung und ihre Folgerungen. Der \emph{eine exakte} Punkt:
$E_\text{bit}=\hbar c/L$ ist eine Photonenenergie.

\textbf{Photon ist Eichboson, Neutrino ist Fermion.} Die Photon-Analogie gilt auf der Ebene
masselos/Dispersion, ist keine Identität (Spin 1 vs.\ $\tfrac12$); der Bezug zur Verbindung gilt dem Eichfeld $A_\mu$ (das Photon ist dessen Quantum), nur im
Photon-Sektor. Die Photon-Phase (Polarisation/Eich-Holonomie) ist dieselbe \emph{Kategorie} wie $\theta$,
nicht dieselbe Zahl. Das Photon sitzt ausserhalb des Leptonen-Massenoperators ($m=0$-Referenz).

\textbf{$\log_2(1/\xi)\approx12{,}87$ ist keine magische Zahl.} Der Gehalt ist Endlichkeit und
$\xi$-Verankerung, nicht der Wert; die genaue Form ($1/\xi$ vs.\ $2/\xi$, Richtungs-/Flavor-Faktor) hängt
an der Modenzählung. Richtung: Geometrie $\to$ Information, nicht umgekehrt.

\textbf{Energieseite: $\xi$-verankert, der Wert-Abgleich ist empirisch.} Theoretisch
$\xi$-/geometrie-festgelegt ($E_0=7{,}398$~MeV, Dok.~011); für die SI-Zahlen der empirische Wert
$\sqrt{m_e m_\mu}=7{,}348$~MeV, damit der Vergleich mit gemessenen Größen passt -- daraus
$m_0,T_0,\tilde T,1/\tilde T$ exakt ($\hbar,c,k_B$ nur Buchhaltung, keine Abweichung). Die einzige
Differenz ist der dimensionslose Theorie-gegen-Empirie-Abgleich $\alpha/\xi$ vs.\ $m_e m_\mu$
($+0{,}68\%$), keine Umrechnung.

\textbf{Bit-Thermodynamik: gerechnet, über die eigene Skala $T_0$.} Das Bit hat die volle
Landauer-Rechnung -- Löschkosten $Q=k_B T\ln2$, relativ zur Eigenenergie $Q/E_0=(T/T_0)\ln2$ (§7); bei
$T=T_0$ ist $Q=E_0\ln2=0{,}69\,E_0$, die intrinsische Größe ist die Bit-Entropie $\ln2$. Keine CMB nötig.
Der UIFT-Parameter $\xi_\text{UIFT}=k_B T_\text{CMB}\ln2/(m_e c^2)\approx3\cdot10^{-10}$ (Dok.~013/243)
ist dieselbe Landauer-Form mit fremden kosmologischen Skalen ($T_\text{CMB}$ statt $T_0$, $m_e$ statt
$E_0$) -- eine andere Größe, nicht FFGFTs geometrisches $\xi=1{,}3\cdot10^{-4}$. Andere Skalenwahl, kein
Scheitern.

\textbf{Neutrino-Argument: spekulativ.} Es ruht weiter auf der spekulativen Entartung (Dok.~007) -- trägt
die These aber nicht mehr, das Photon trägt sie.

\section{Bilanz: nicht ob, sondern welcher Kanal}
Die Betrag/Phase-Arbeit verschiebt die Informationsfrage von \enquote{ist Information fundamental?} zu
\enquote{welcher Kanal?}. Vopson/Landauer sind der Betragskanal -- real, aus der $T^4$-Geometrie
hergeleitet, aber blind für die Phase. Der Phasenkanal (Mischung, Holonomie, $\theta$) ist unabhängige,
relationale Information; in der Signalverarbeitung trägt gerade er die Struktur. Der reinste Anker ist das
Photon: exakt masselos, FFGFTs Informationseinheit $E_\text{bit}=\hbar c/L$ selbst, und das Quantum der
$U(1)$-Verbindung -- begrifflich \enquote{it from connection}. Die \emph{Menge} an Information ist
absolut und in $\xi$ verankert ($I_1=\log_2(1/\xi)\approx12{,}87$ Bit); Energie und Kapazität sind ihre
dimensionalen Beiwerke. Nur die Geometrie hält beide Kanäle -- ehrlich abgegrenzt als Diagnose, nicht als
neue Herleitung.

